Start with an i.i.d sequence $X_1, X_2, \ldots$ with moment condition $E e^{t X_1} < \infty$ for all $t \in \mathbb{R}$. We also assume for now that
\begin{enumerate}
    \item $E X_1 = 0$. This can be easily removed by centering the random variable. $P\left( \frac{S_n}{n} \ge a \right) = P\left( \frac{S_n - n \mu}{n} a - \mu \right) $ where $\varphi_{X - \mu}(t) = \varphi_X(t) e^{- \mu t} < \infty $.
    \item $X_1$ has both positive and negative tails. The one-sided case can be handled very simply. Suppose that $P(X_1 \le 0) = 1$. If $P(X_1 = 0) = \rho > 0$, we have
	    \[
	    	P\left( \frac{S_n}{n} \ge 0 \right)  = P(X_1 = 0)^n = \rho^n
	    .\] 
	    On the other hand, if $P(X_1 < 0) = 1$, we have
	     \[
	    	P\left( \frac{S_n}{n}\ge 0 \right)  = 0
	    .\] 
	    In both cases the result follows.
\end{enumerate}
We want to obtain the large deviation principle for $S_n / n$, where $S_n = X_1 + \ldots + X_n$. In other words, we want to obtain information about exponential decay of the probability
\[
    P\left( \frac{S_n}{n} > x \right)
\]
We may still simplify the proof by considering 
\begin{enumerate}[resume]
    \item $x = 0$. Because $P(\frac{S_n}{n}>x) = P\left( \frac{S_n - n x }{n} > 0 \right) $ which is treated by assumption (i).
\end{enumerate}

We start by applying an exponential change of measure:
\begin{align*}
    P\left( \frac{S_n}{n} > 0 \right) &= P\left( S_n > 0 \right) \\
    &= E\left[ 1_{S_n > 0} \right] \\
    &= E\left[ 1_{S_n > 0} \frac{e^{\lambda S_n}}{ E \left[ e^{\lambda S_n} \right]} e^{-\lambda S_n} \right] E \left[ e^{\lambda S_n} \right] \\
    &= \tilde E \left[ 1_{\hat S_n > 0} e^{-\lambda \hat S_n} \right] E \left[ e^{\lambda S_n} \right] \\
\end{align*}

We use the notation $\tilde E$ to denote the expectation with respect to the new measure. We denote by $\hat S_n$ the random variable $S_n$ under the new measure. 

The second term $E \left[ e^{\lambda S_n} \right]$ contains the information about the tails of the distribution: $\frac{1}{n} \log E \left[ e^{\lambda S_n} \right] =: \Lambda_n(\lambda) \to \Lambda(\lambda)$. We want to find the $\lambda$ such that this term contains \textit{all the information} about the tails of the distribution, i.e. the other term is $e^{- o(n)}$. It turns out that there is only one such $\lambda^*$, and we say $I(0) = -\Lambda(\lambda^*)$, giving the LDP 
\[
    \frac{1}{n} \log P\left( \frac{S_n}{n} > 0 \right) = - I(0) + o(1).
\]

Look at the left hand side. One notices that the term $e^{-\lambda \hat S_n}$ grows or decays exponentially fast unless $\hat S_n$ is centered, in which case it is $O(\sqrt{n})$ due to the central limit theorem (Note that $\hat S_n$ inherits all moment conditions, and all summands remain independent due to the product structure of the Radon-Nikodym derivative.) How can we center $\hat S_n$? Let's compute the moment generating function of $\hat S_n$:
\begin{align*}
    \tilde \varphi_{\hat S_n}(t) &= \tilde E \left[ e^{t \hat S_n} \right] \\
    &= E \left[ e^{t S_n} \frac{e^{\lambda S_n}}{E \left[ e^{\lambda S_n} \right]} \right] \\
    &= \varphi_{S_n}(t + \lambda) / \varphi_{S_n}(\lambda)
\end{align*}

If we take $\lambda$ such that $\varphi'(\lambda) = 0$, we have
\[
    \tilde E[\hat S_n] = \tilde \varphi_{\hat S_n}'(0) = \frac{\varphi_{S_n}'(\lambda)}{\varphi_{S_n}(\lambda)} = 0.
\]

To finish the proof one performs the following sanity checks:
\begin{tcolorbox}[]
    \item When $X_1$ has both positive and negative tails, $\varphi_{S_n}(\lambda)$ is continuous, differentiable, strictly convex and has a unique minimum at $\lambda^* \in \mathbb{R}$. 
\end{tcolorbox}

The proof is already complete, but just for demonstration, we give the proof for $x \neq 0$. 
\begin{align*}
    P\left( S_n > n x \right)
    &= E\left[ 1_{S_n > n x} \right] \\
    &= E\left[ 1_{S_n > n x} \frac{e^{\lambda S_n - n \lambda x}}{ E \left[ e^{\lambda S_n - n \lambda x} \right]} e^{-\lambda S_n + n \lambda x} \right] E \left[ e^{\lambda S_n - n \lambda x} \right] \\
    &= \tilde E \left[ 1_{\hat S_n > n x} e^{-\lambda \hat S_n + n \lambda x} \right] E \left[ e^{\lambda S_n - n \lambda x} \right] \\
\end{align*}

First look at the second term. We have
\[
    \frac{1}{n} \log E \left[ e^{\lambda S_n - n \lambda x} \right] = \frac{1}{n} \log E \left[ e^{\lambda S_n} \right] - \lambda x = \Lambda_n(\lambda) - \lambda x
    \to \Lambda(\lambda) - \lambda x.
\]

% If we minimize over $\lambda$, we can have the exponential decay: $-I(x) = \inf_{\lambda} \left( \Lambda(\lambda) - \lambda x \right)$.
It turns out that the minimizer $\lambda^*$ is exactly the one making the first term $e^{o(n)}$. By the same argument as above, we can center $\hat S_n - n x$ by taking $\lambda$ such that $\varphi_{S_n - n x}'(\lambda) = 0$. 
We have
\[
    E [ (S_n - n x) e^{\lambda S_n - n \lambda x} ] = 0
,\]
i.e., 
\[
    E [S_n e^{\lambda S_n} ] = n x E [ e^{\lambda S_n} ],
\]
i.e.,
\[
    \varphi_{S_n}'(\lambda) = n x \varphi_{S_n}(\lambda),
\]
i.e.,
\[
    \Lambda_n'(\lambda) = x.
\]

When $\Lambda_n(\lambda)$ is differentiable and strictly convex in some domain, such $\lambda$ is the unique maximizer in the Legendre transform
\[
    I(x) = \sup_{\lambda} \left( \lambda x - \Lambda(\lambda) \right).
\]

Relaxing assumptions (1) and (2) are left as an exercise. Let's state the result:

\begin{theorem}[Cramer's theorem]
    Let $X_1, X_2, \ldots$ be an i.i.d. sequence of random variables with moment condition $E e^{t X_1} < \infty$ for all $t \in \mathbb{R}$. Then the sequence $S_n = X_1 + \ldots + X_n$ satisfies the large deviation principle with rate function
    \[
        I(x) = \sup_{\lambda} \left( \lambda x - \Lambda(\lambda) \right).
    \]
\end{theorem}

For properties of the rate function, see \autocite[Lemma~I.14]{denhollanderLargeDeviations2008}. 


